\subsection{Cause 5: Type-Unaware Variable Naming in Bottom-Clause Generalization}
\label{sec:andante-cause5}

A fifth defect was identified in Andante's construction of the generalized
bottom clause (Algorithm~1 of \citet{muggleton1995inverse}). During
bottom-clause construction, every ground constant bound to a $+$-typed
(input) or $-$-typed (output) mode argument must be replaced by a fresh
program variable, so that the resulting clause generalizes beyond the
single seed example it was built from. This renaming was performed by
\texttt{Constant.to\_variable\_name()}, which derived the variable's
display symbol \emph{solely} from the constant's raw value (e.g.\ the
integer $1$ always produced the symbol \texttt{V1}), with no reference to
the constant's declared Progol type.

This is unsound whenever two \emph{different}-typed constants occurring
as distinct arguments of the \emph{same} literal happen to share the same
raw value. In the \textsc{iggp-rps} (general game playing, rock-paper-scissors)
domain, the target example
$\mathtt{next\_score(1, p2, 1)}$ binds the game-state identifier
(declared type \texttt{ex}) and the numeric score (declared type
\texttt{int}) to the same value $1$. Because naming ignored type, both
occurrences were mapped to the identical variable \texttt{V1}, and the
generalized head atom collapsed from the intended
$\mathtt{next\_score(V_{ex}, P, V_{int})}$ to the malformed
$\mathtt{next\_score(V1, P2, V1)}$ — a single variable erroneously
occupying two semantically unrelated argument positions. This
ill-formed clause later caused unification to attempt binding one
variable to two incompatible types in sequence, raising an unhandled
\texttt{AttributeError} deep inside the substitution engine and aborting
the learning process entirely.

\paragraph{Fix.} \texttt{to\_variable\_name()} was extended with an
optional \texttt{type\_name} parameter (default \texttt{None}, preserving
the method's original zero-argument contract). Within
\texttt{build\_bottom\_i}, a local collision table records, for each
literal being generalized, the type first associated with a given base
variable name; if a \emph{second}, differently-typed constant would map
to that same base name \emph{within the same literal}, it is instead
assigned a disambiguated name (e.g.\ \texttt{V1\_int}). A second table,
\texttt{head\_renames}, records the names chosen while generalizing the
head atom and is consulted while generalizing every subsequent body
literal, so that a value disambiguated in the head (such as the
\texttt{int}-typed $1$) is named identically wherever it recurs in the
clause (for instance, as the output of \texttt{my\_succ/2}). Without this
propagation, the head and the relevant body literal would use different
symbols for what must be the same variable, and the correct clause would
remain syntactically unreachable by the refinement search.

\paragraph{Scope of the fix.} The collision check is deliberately local
to a single literal's own argument list, rather than applied to every
recurrence of a value throughout the whole bottom clause. An earlier,
more aggressive version of the fix disambiguated \emph{any} same-valued,
differently-typed constant pair anywhere in the clause; while this also
corrected \textsc{iggp-rps}, it substantially degraded search performance
on the Zendo dataset (from $\approx$18\,s to over 170\,s without
convergence), because Zendo's state, piece, and attribute identifiers are
frequently small integers that legitimately co-occur across different
literals and are intentionally merged into shared variables by Andante's
anti-unification mechanism. Restricting disambiguation to genuine
intra-literal collisions preserves this beneficial sharing everywhere it
is safe, and only intervenes in the narrow case that is actually
unsound.

\paragraph{Validation.} With this fix, \textsc{iggp-rps} converges to
three clauses structurally identical to the reference solution obtained
independently by Popper, achieving perfect coverage
($tp=108$, $fn=0$, $tn=356$, $fp=0$; accuracy, precision, recall and
$F_1$ all equal to $1.0$) in approximately 66--90\,s. No regression was
observed on any previously-validated dataset (\texttt{family},
\texttt{short\_family}, \texttt{trains}, \texttt{trains2}, \texttt{zendo},
\texttt{zendo1}), each of which continues to report identical hypotheses,
accuracy and running time to those measured prior to this fix.

\begin{table}[h]
\centering
\begin{tabular}{@{}lll@{}}
\toprule
Dataset & Accuracy / F$_1$ & Runtime \\
\midrule
\textsc{iggp-rps} (after fix) & 1.000 / 1.000 & $\approx$66--90\,s \\
\textsc{iggp-rps} (before fix) & crash (\texttt{AttributeError}) & --- \\
\bottomrule
\end{tabular}
\caption{Effect of the Cause 5 fix on the \textsc{iggp-rps} dataset.}
\end{table}
